\section{Method}\label{sec:method}
We now describe GCD and how it can constrain the output of LMs at decoding time based on a formal grammar.
%
We first explain how to specify the output spaces of various NLP tasks via formal grammars.
Then we show how an incremental parser can be used as a completion engine to constrain the LM's generation process to produce grammatically valid outputs only.

\subsection{NLP tasks as formal languages}\label{subsec:define-the-output-of-nlp-tasks-as-formal-languages}

The input $x$ and output $y$ of NLP tasks are typically sequences of tokens, $x = \langle x_0, \dots, x_{n-1} \rangle$ and $y = \langle y_0, \dots, y_{m-1} \rangle$.
Whereas the input $x$ is usually arbitrary, for many tasks the output $y$ needs to follow a specific structure.
For instance, in information extraction, $y$ is required to consist of subject--relation--object triplets (\cf\ \Figref{fig:intro_figure}).
Since formal languages provide a rigorous and complete framework for describing the structure of any computable set of object (according to the Church--Turing thesis), they offer a promising way to define the output spaces of structured NLP tasks.
In order to define the formal languages corresponding to the output spaces of structured NLP tasks, our framework relies on formal grammars, a universal formalism that can describe any formal language.
For tractability, we focus on the class of context-free grammars.

% tractable formal grammar\footnote{Regular or context-free grammars are considered tractable.} as demonstrated in \Figref{fig:intro_figure} and \ref{fig:grammars_showcase_more}.


% Formal languages provide a rigorous and complete framework for describing the structure of any computable object (according to the Church--Turing thesis).
% The output of an NLP task can be formalized as a tuple  ($X$, $Y$), where $X$ is the input sequence and $Y$ is the output sequence.
% Both the input and output are typically sequences of tokens, $x = \langle x_1, x_2, ..., x_n \rangle$ and $y = \langle y_1, y_2, ..., y_m \rangle$.
% While the input sequence is usually arbitrary, the output sequence can, for many tasks, be restricted to a specific structure.
% Interestingly, many NLP tasks involve structured outputs that can be effectively described using a tractable formal grammar\footnote{Regular or context-free grammars are considered tractable.} as demonstrated in \Figref{fig:intro_figure} and \ref{fig:grammars_showcase_more}.

\xhdr{Token-level formal grammars}
A formal grammar $G$ is defined as a tuple $(V, \Sigma, P, S)$ where
\begin{itemize}
\denselist
\item $V$ is a finite set of non-terminal symbols,
\item $\Sigma$ is a finite set of terminal symbols,
\item $P$ is a finite set of production rules,
\item $S \in V$ is the start symbol.
\end{itemize}

We illustrate the suitability of formal grammars for specifying the output spaces of structured NLP tasks in \Figref{fig:grammars_showcase_more} using 14 common tasks as examples.

In our approach, the user first writes a formal grammar $G$ over characters.
In order to obtain a token-level grammar $G_\text{tok}$ that can be used to constrain the direct output of LMs---token sequences---we use the token set $\Sigma_\text{tok}$ as terminal symbols and
apply the tokenizer to the sequences of terminal symbols appearing in the rules $P$, obtaining the token-level rules $P_\text{tok}$.
This yields the token-level grammar $G_\text{tok} = (V, \Sigma_\text{tok}, P_\text{tok}, S)$, which describes the same language as the character-level grammar $G$.
We then use an incremental parser (see below) to decide whether a token sequence $y$ is in the language generated by $G_\text{tok}$.

Despite its simplicity, this approach has some limitations.
Widely used tokenization methods such as BPE~\citep{bpe-sennrich-etal-2016-neural} allow the same string to be tokenized in different ways.
For example, the string ``\texttt{[[[}'' can be tokenized as ``\texttt{[[ [}'', ``\texttt{[[[}'' or ``\texttt{[ [ [}'',
all of which would be detokenized to the original string ``\texttt{[[[}''.
%
To avoid this ambiguity, we can add extra spaces between the brackets in the grammar and force the single bracket to be a token.
This approach is, however, not principled and relies on a specific tokenizer.
We therefore believe that a principled approach for defining token-level grammars is an interesting direction for future work.


\xhdr{Grammatical Framework}
Since the token-level grammar $G\textsubscript{tok}$ is tokenizer-dependent, there are in general multiple token-level grammars for the same grammar $G$.
This one-to-many mapping from a character-level grammar to a token-level grammar is analogous to the one-to-many mapping from abstract syntax trees to different programming languages.
For this reason, we adopt Grammatical Framework (GF)~\cite{ranta-2019-grammatical} to define both the grammar $G$ and the token-level grammar $G_\text{tok}$.
GF is a meta-language for multilingual grammar applications, which allows us to define an abstract grammar as well as concrete grammars for linearizing abstract syntax trees into different tokenizer\hyp specific ``languages''.
In our case, the abstract grammar is the character-level grammar $G$ and the concrete grammars are the token-level grammars for different tokenizers.


\xhdr{Input-dependent grammars (IDG)}
While the output of many NLP tasks can be described by a formal grammar, some tasks require a grammar that is dependent on the input sequence.
For example, in entity dismbiguation, the output needs to be constrained to the set of entity candidates, which usually contains dozens of entity names semantically related to the target mention in the input sequence.
In constituency parsing, outputs are parse trees whose terminal nodes are the input tokens.
These two tasks both require a grammar that is dependent on the input.
%
Existing work has focused on using a single grammar to constrain the decoding regardless of the input sequence.
Whereas this is suitable for tasks where the output space is independent of the input sequence, such as code generation~\cite{poesia2022synchromesh} or information extraction~\cite{josifoski-etal-2022-genie} (\cf\ \Figref{fig:intro_figure}),
13 of the 14 tasks listed in \Figref{fig:grammars_showcase_more} (those with an asterisk) require an input-dependent grammar.


\subsection{Grammar-constrained decoding (GCD)}\label{subsec:grammar-constrained-decoding}

The LM decoding process produces tokens one by one.
To enforce the formal grammar, we intervene during decoding by pruning the probability distribution to include only the subset of tokens that are allowed by the formal grammar.
The subset of allowed tokens is returned by an incremental parser, which takes the partially generated sequence and the formal grammar as inputs and returns the set of next allowed tokens.
The role of the parser can be abstracted as a \textit{completion engine}~\cite{poesia2022synchromesh}, which derives completion candidates from the partial sequence and the formal grammar $G$.
In this work, we use the \textit{incremental parser} of Grammatical Framework~\cite{angelov-2009-incremental} as the completion engine.
This process is compatible with any decoding algorithm, including greedy decoding, beam search, top-$k$ sampling, etc.
GCD can also be applied to any autoregressive language model, provided that we have access to the distribution over the vocabulary at each decoding step.
Since API-based services such as OpenAI do not provide access to the distribution over the vocabulary, they cannot be used with GCD.

\subsection{Few-shot learning with GCD}\label{subsec:few-shot-learning-with-gcd}

To adapt a pretrained LM to a new task, one can either finetune the LM on task-specific training data or use few-shot learning if the LM is powerful enough.
In this work, we use GCD in conjunction with few-shot learning to adapt a pretrained large LM (LLM) to new tasks.
Instead of prompting the LLM to generate free-form text, we use GCD to constrain the output to be grammatically correct.
This allows us to leverage the LLM's few-shot learning capability together with the grammar-induced knowledge of the task's output structure.

For the same task, we can use different grammars to constrain the output of the LLM.
For example, in entity disambiguation, we can use a grammar $G_1$ that depends on the input sequence to constrain the output to the input-specific candidate set, or we could use a grammar $G_2$ that is independent of the input to constrain the output to be any valid entity name.
While both $G_1$ and $G_2$ can be used to constrain the output of the LLM, $G_1$ reduces the search space more and thus is more effective.

% In some cases, we cannot find a perfect grammar to describe the output of the task, \eg, constituency parsing, where the whole task is to approximate the underlying constituency grammar.
% A relaxed grammar can still provide useful guidance to the LM and reduce errors.


